% !TeX root = closed-systems-from-comparison-completeness.tex
% ============================================================
\chapter{Descent--Obstruction Classification}
\label{ch:descent_obstruction}
% ============================================================
\begin{chapterorientation}
\orientationitem{Input}{The finite triangle witness and a connected quotient protocol with representative lifts.}
\orientationitem{Role}{This chapter globalizes finite transport failure into a loop-representation classification.}
\orientationitem{Output}{Tree normalization, loop data, and endpoint-determined transport criteria for the rigidity bridge.}
\end{chapterorientation}
\section{Introduction}
Under the standing principle of closed-world admissibility
(\cref{sp:closed-world}), this chapter globalizes the transport-visibility
clause \cref{sp:transport} from finite triangle witnesses to loop-level
obstruction representations.
Starting from the first finite witness isolated in \cref{ch:triangles},
it develops the global obstruction classification via loop
representations required in \cref{ch:bridge}.
\Cref{ch:triangles} identified the first finite witness of failure of a
common diagonal gauge (\cref{thm:triangle}). The present chapter addresses the complementary
global problem: how does one classify all representative lifts of a
connected quotient protocol up to gauge?
The answer is controlled by a tree--loop decomposition.
Fix a spanning tree in the underlying undirected graph of
\(\mathcal I\). After a suitable gauge
normalization, all representatives become constant and every morphism in
the tree groupoid carries trivial transport. All residual transport then
lives on loops. More precisely, the remaining datum is a homomorphism
from the based loop group of the groupoid completion
\[
	\Omega_{i_0}(\mathcal G(\mathcal I))
\]
into the stabilizer subgroup of a chosen base representative.
This loop representation is the intrinsic obstruction to descent.
Tree transport carries no gauge-invariant content; loop transport is the
residue. Endpoint-determined transport on the groupoid completion is
equivalent to triviality of this residue.
The chapter proves three things.
\begin{enumerate}
	\item Every lift is gauge-equivalent to a tree-normalized lift with
	      constant representative choice.
	\item Tree-normalized lifts are completely determined by their based loop
	      representation.
	\item Gauge-equivalence classes of lifts are classified by conjugacy
	      classes of homomorphisms
	      \[
		      \rho:\Omega_{i_0}(\mathcal G(\mathcal I))\to H_*,
		      \qquad
		      H_*=\Stab(x_*),
	      \]
	      where \(x_*\) is the chosen base representative.
\end{enumerate}
Thus the morphism locus of the two-locus principle admits a complete
and explicit obstruction theory. Beyond the already-separated
object-level representative data, no additional morphism-level
obstruction parameter appears.
% ============================================================
\section{Connected quotient protocols and their lifts}
% ============================================================
Let
\[
	\overline R:\mathcal I\to \Phys^{\disc}
\]
be a quotient-level protocol, where \(\mathcal I\) is a connected small
category and \(\Phys^{\disc}\) is the discrete category on the quotient
state space \(\Phys\).
By \cref{lem:lift-existence}, connectedness of \(\mathcal I\) implies
that \(\overline R\) is constant. We therefore fix the unique point
\[
	p\in \Phys
\]
such that
\[
	\overline R(i)=p
	\qquad
	\text{for all }i\in \Obj(\mathcal I).
\]
Choose once and for all a representative
\[
	x_*\in X
	\qquad
	\text{with}\qquad
	\pi(x_*)=p,
\]
and set
\[
	H_*:=\Stab(x_*)
	=
	\{h\in G:\ h\cdot x_*=x_*\}.
\]
We work with representative lift data \((x,g)\) in the sense of
\cref{thm:lift-classification-main}: thus
\[
	x_i\in X
	\qquad (i\in \Obj(\mathcal I)),
\]
and
\[
	g_f\in G
	\qquad (f\in \Mor(\mathcal I)),
\]
satisfy
\begin{align}
	g_f\cdot x_i  & = x_j
	              &             & \text{for every }f:i\to j, \label{eq:DO-T1}                     \\
	g_{\id_i}     & = e
	              &             & \text{for every }i, \label{eq:DO-T2}                            \\
	g_{f'\circ f} & = g_{f'}g_f
	              &             & \text{for every composable }i\xrightarrow{f}j\xrightarrow{f'}k.
	\label{eq:DO-T3}
\end{align}
Equivalently, these data determine a functor
\[
	T:\mathcal I\to X\sslash G
\]
lifting \(\overline R\).
% ============================================================
\section{Tree normalization}
% ============================================================
We now isolate the gauge-removable part of the transport.
\begin{definition}[Underlying graph and spanning tree]
	Let \(|\mathcal I|\) denote the underlying undirected graph of
	\(\mathcal I\): its vertices are the objects of \(\mathcal I\), and an
	undirected edge joins \(i\) and \(j\) whenever \(\mathcal I\) contains a
	morphism \(i\to j\) or \(j\to i\).
	A \emph{spanning tree} \(\mathcal T\subseteq |\mathcal I|\) is a
	connected, acyclic spanning subgraph.
\end{definition}
Fix a base object
\[
	i_0\in \Obj(\mathcal I)
\]
and a spanning tree \(\mathcal T\subseteq |\mathcal I|\).
Let
\[
	\mathcal G(\mathcal I)
\]
denote the groupoid completion of \(\mathcal I\), as in
\cref{def:groupoid-completion}. Let
\[
	\mathcal G(\mathcal T)\subseteq \mathcal G(\mathcal I)
\]
denote the subgroupoid generated by the tree edges. Since \(\mathcal T\)
is a tree, for each object \(i\) there is a unique morphism
\[
	\gamma_i:i_0\to i
\]
in \(\mathcal G(\mathcal T)\).
\begin{lemma}[Tree-path transport]
	\label{lem:DO-tree-path}
	For every object \(i\),
	\[
		g_{\gamma_i}\cdot x_{i_0}=x_i.
	\]
\end{lemma}
\begin{proof}
	By \cref{lem:cocycle-extension}, the transport cocycle extends uniquely
	from \(\mathcal I\) to a functor
	\[
		\widetilde T:\mathcal G(\mathcal I)\to X\sslash G.
	\]
	Write
	\[
		\widetilde T(\gamma_i)=(g_{\gamma_i},x_{i_0}).
	\]
	Since \(\gamma_i:i_0\to i\), the target of this morphism is \(x_i\).
	By definition of the action groupoid,
	\[
		g_{\gamma_i}\cdot x_{i_0}=x_i.
	\]
\end{proof}
\begin{theorem}[Tree normalization]
	\label{thm:DO-tree-normalize}
	There exists a gauge modification \(h=(h_i)_{i\in\Obj(\mathcal I)}\)
	such that the gauge-modified lift datum \((x^h,g^h)\) satisfies:
	\begin{enumerate}
		\item
		      \[
			      x_i^h=x_*
			      \qquad
			      \text{for all }i\in\Obj(\mathcal I);
		      \]
		\item for every morphism \(\tau\) in the tree groupoid
		      \(\mathcal G(\mathcal T)\),
		      \[
			      g_\tau^h=e.
		      \]
	\end{enumerate}
	In particular, every tree edge carries trivial transport after
	normalization.
\end{theorem}
\begin{proof}
	Because \(\overline R\) is constant with value \(p\), every representative
	\(x_i\) lies in the orbit \(\pi^{-1}(p)\). In particular,
	\[
		\pi(x_{i_0})=\pi(x_*)=p.
	\]
	Hence there exists \(u\in G\) such that
	\[
		u\cdot x_{i_0}=x_*.
	\]
	By \cref{lem:DO-tree-path},
	\[
		g_{\gamma_i}\cdot x_{i_0}=x_i
		\qquad
		\text{for all }i.
	\]
	Define
	\[
		h_i:=u\,g_{\gamma_i}^{-1}
		\qquad
		(i\in\Obj(\mathcal I)).
	\]
	We first compute the transformed representatives. By definition of gauge
	modification,
	\[
		x_i^h=h_i\cdot x_i
		=
		u\,g_{\gamma_i}^{-1}\cdot x_i.
	\]
	Using \(x_i=g_{\gamma_i}\cdot x_{i_0}\), we obtain
	\[
		x_i^h
		=
		u\,g_{\gamma_i}^{-1}\cdot(g_{\gamma_i}\cdot x_{i_0})
		=
		u\cdot x_{i_0}
		=
		x_*.
	\]
	Thus all transformed representatives are equal to the base
	representative \(x_*\).
	Next let \(\tau:i\to j\) be any morphism in the tree groupoid
	\(\mathcal G(\mathcal T)\). Then in \(\mathcal G(\mathcal T)\) one has
	\[
		\tau=\gamma_j\circ \gamma_i^{-1}.
	\]
	Hence, by functoriality of the extended transport,
	\[
		g_\tau=g_{\gamma_j}g_{\gamma_i}^{-1}.
	\]
	The transformed transport label is therefore
	\[
		g_\tau^h
		=
		h_j g_\tau h_i^{-1}
		=
		(u g_{\gamma_j}^{-1})(g_{\gamma_j}g_{\gamma_i}^{-1})(g_{\gamma_i}u^{-1})
		=
		e.
	\]
	This proves the second statement.
\end{proof}
\begin{definition}[Tree-normalized lift]
	A representative lift datum \((x,g)\) is called
	\emph{tree-normalized} if
	\[
		x_i=x_*
		\qquad
		\text{for all }i\in\Obj(\mathcal I),
	\]
	and
	\[
		g_\tau=e
		\qquad
		\text{for all }\tau\in \Mor(\mathcal G(\mathcal T)).
	\]
\end{definition}
\begin{remark}[What remains after normalization]
	After tree normalization, the object-level data are constant and the
	entire tree transport vanishes. All residual information is therefore
	concentrated on loops. This is the descent--obstruction form of the
	morphism locus.
\end{remark}
% ============================================================
\section{Loop representations}
% ============================================================
Let
\[
	\Omega_{i_0}(\mathcal G(\mathcal I))
	:=
	\Hom_{\mathcal G(\mathcal I)}(i_0,i_0)
\]
denote the based loop group at \(i_0\), with group law given by
composition in \(\mathcal G(\mathcal I)\).
\begin{lemma}[Loop holonomy lands in the stabilizer]
	\label{lem:DO-holonomy}
	If \((x,g)\) is tree-normalized, then for every loop
	\[
		\ell\in \Omega_{i_0}(\mathcal G(\mathcal I))
	\]
	one has
	\[
		g_\ell\in H_*=\Stab(x_*).
	\]
\end{lemma}
\begin{proof}
	Since \((x,g)\) is tree-normalized, the representative at the base object
	is \(x_*\). Because \(\ell\) is a loop at \(i_0\), the corresponding
	morphism in the action groupoid is
	\[
		(g_\ell,x_*):x_*\to x_*.
	\]
	Thus
	\[
		g_\ell\cdot x_*=x_*,
	\]
	which is exactly the statement that \(g_\ell\in H_*\).
\end{proof}
\begin{definition}[Obstruction representation]
	\label{def:DO-rho}
	For a tree-normalized lift datum \((x,g)\), define
	\[
		\rho_g:\Omega_{i_0}(\mathcal G(\mathcal I))\to H_*,
		\qquad
		\rho_g(\ell):=g_\ell.
	\]
\end{definition}
\begin{lemma}
	\label{lem:DO-rho-hom}
	The map \(\rho_g\) is a group homomorphism.
\end{lemma}
\begin{proof}
	Let \(\ell_1,\ell_2\in \Omega_{i_0}(\mathcal G(\mathcal I))\). Since
	transport is functorial on the groupoid completion,
	\[
		g_{\ell_2\circ \ell_1}=g_{\ell_2}g_{\ell_1}.
	\]
	Therefore
	\[
		\rho_g(\ell_2\circ \ell_1)
		=
		g_{\ell_2\circ \ell_1}
		=
		g_{\ell_2}g_{\ell_1}
		=
		\rho_g(\ell_2)\rho_g(\ell_1).
	\]
	Thus \(\rho_g\) is a homomorphism.
\end{proof}
\begin{lemma}[Gauge conjugation of the loop representation]
	\label{lem:DO-rho-conj}
	Let \((x,g)\) be tree-normalized, and let \(k\in H_*\). Define a gauge
	transformation by
	\[
		h_i:=k
		\qquad
		\text{for all }i\in\Obj(\mathcal I).
	\]
	Then the transformed tree-normalized lift datum satisfies
	\[
		\rho_{g^h}(\ell)=k\,\rho_g(\ell)\,k^{-1}
		\qquad
		\text{for all }\ell\in \Omega_{i_0}(\mathcal G(\mathcal I)).
	\]
\end{lemma}
\begin{proof}
	For any morphism \(\alpha:i\to j\) in \(\mathcal G(\mathcal I)\), the
	gauge transformation rule gives
	\[
		g_\alpha^h=h_j g_\alpha h_i^{-1}=k g_\alpha k^{-1}.
	\]
	Applying this to a loop \(\ell:i_0\to i_0\) yields
	\[
		\rho_{g^h}(\ell)
		=
		g_\ell^h
		=
		k g_\ell k^{-1}
		=
		k\,\rho_g(\ell)\,k^{-1}.
	\]
\end{proof}
% ============================================================
\section{Reconstruction from the loop representation}
% ============================================================
We now prove the converse statement: a homomorphism into the base
stabilizer determines a unique tree-normalized lift.
\begin{definition}[Loop reduction map]
	For any morphism
	\[
		\alpha:i\to j
	\]
	in \(\mathcal G(\mathcal I)\), define the associated based loop
	\[
		\lambda(\alpha):=\gamma_j^{-1}\circ \alpha\circ \gamma_i
		\in
		\Omega_{i_0}(\mathcal G(\mathcal I)).
	\]
\end{definition}
\begin{lemma}[Multiplicativity of loop reduction]
	\label{lem:DO-lambda}
	For composable morphisms
	\[
		i\xrightarrow{\alpha} j\xrightarrow{\beta} k
	\]
	in \(\mathcal G(\mathcal I)\),
	\[
		\lambda(\beta\circ \alpha)=\lambda(\beta)\circ \lambda(\alpha).
	\]
\end{lemma}
\begin{proof}
	By definition,
	\[
		\lambda(\beta\circ \alpha)
		=
		\gamma_k^{-1}\circ \beta\circ \alpha\circ \gamma_i.
	\]
	Insert the identity
	\[
		\id_j=\gamma_j\circ\gamma_j^{-1}
	\]
	between \(\beta\) and \(\alpha\). Then
	\[
		\lambda(\beta\circ \alpha)
		=
		(\gamma_k^{-1}\circ \beta\circ \gamma_j)
		\circ
		(\gamma_j^{-1}\circ \alpha\circ \gamma_i)
		=
		\lambda(\beta)\circ \lambda(\alpha).
	\]
\end{proof}
\begin{proposition}[Reconstruction from a loop representation]
	\label{prop:DO-reconstruct}
	Let
	\[
		\rho:\Omega_{i_0}(\mathcal G(\mathcal I))\to H_*
	\]
	be a homomorphism. Then there exists a unique tree-normalized
	representative lift datum \((x,g)\) such that
	\[
		\rho_g=\rho.
	\]
\end{proposition}
\begin{proof}
	Define
	\[
		x_i:=x_*
		\qquad
		\text{for all }i\in\Obj(\mathcal I).
	\]
	For any morphism \(\alpha:i\to j\) in \(\mathcal G(\mathcal I)\), define
	\[
		g_\alpha:=\rho(\lambda(\alpha))\in H_*\subseteq G.
	\]
	We first check that this determines a functor
	\[
		\widetilde T_\rho:\mathcal G(\mathcal I)\to X\sslash G.
	\]
	Because \(g_\alpha\in H_*=\Stab(x_*)\), one has
	\[
		g_\alpha\cdot x_i=g_\alpha\cdot x_*=x_*=x_j.
	\]
	Hence
	\[
		(g_\alpha,x_*):x_i\to x_j
	\]
	is a well-defined morphism in the action groupoid.
	For the identity morphism \(\id_i\), one has
	\[
		\lambda(\id_i)=\gamma_i^{-1}\circ \id_i\circ \gamma_i=\id_{i_0},
	\]
	so
	\[
		g_{\id_i}=\rho(\id_{i_0})=e.
	\]
	For composable \(\alpha,\beta\), \cref{lem:DO-lambda} and the
	homomorphism property of \(\rho\) give
	\[
		g_{\beta\circ \alpha}
		=
		\rho(\lambda(\beta\circ \alpha))
		=
		\rho(\lambda(\beta)\circ \lambda(\alpha))
		=
		\rho(\lambda(\beta))\rho(\lambda(\alpha))
		=
		g_\beta g_\alpha.
	\]
	Thus \(\widetilde T_\rho\) is a functor.
	Restricting from \(\mathcal G(\mathcal I)\) to \(\mathcal I\) gives a
	representative lift datum \((x,g)\).
	It remains to verify tree-normalization. If \(\tau:i\to j\) is a
	morphism in the tree groupoid \(\mathcal G(\mathcal T)\), then
	\[
		\tau=\gamma_j\circ \gamma_i^{-1},
	\]
	hence
	\[
		\lambda(\tau)
		=
		\gamma_j^{-1}\circ \tau\circ \gamma_i
		=
		\gamma_j^{-1}\circ \gamma_j\circ \gamma_i^{-1}\circ \gamma_i
		=
		\id_{i_0}.
	\]
	Therefore
	\[
		g_\tau=\rho(\id_{i_0})=e.
	\]
	So the lift is tree-normalized.
	Finally, for a loop \(\ell\in \Omega_{i_0}(\mathcal G(\mathcal I))\),
	\[
		\lambda(\ell)
		=
		\gamma_{i_0}^{-1}\circ \ell\circ \gamma_{i_0}
		=
		\ell,
	\]
	because \(\gamma_{i_0}=\id_{i_0}\). Hence
	\[
		\rho_g(\ell)=g_\ell=\rho(\ell).
	\]
	Thus \(\rho_g=\rho\).
	Uniqueness is immediate: a tree-normalized lift is determined by its
	values on all morphisms, and for each \(\alpha\) those values are determined
	to be \(\rho(\lambda(\alpha))\).
\end{proof}
% ============================================================
\section{Descent--obstruction classification}
% ============================================================
We now combine normalization, reconstruction, and gauge conjugation into
the main theorem.
\begin{theorem}[Descent--obstruction classification]
	\label{thm:DO-classification}
	Under the standing principle \cref{sp:closed-world}, let
	\(\mathcal I\) be connected, fix a base object \(i_0\), and fix a
	representative
	\[
		x_*\in \pi^{-1}(p),
		\qquad
		H_*=\Stab(x_*).
	\]
	Then gauge-equivalence classes of representative lifts of the connected
	quotient protocol
	\[
		\overline R:\mathcal I\to \Phys^{\disc}
	\]
	are classified by \(H_*\)-conjugacy classes of homomorphisms
	\[
		\rho:\Omega_{i_0}(\mathcal G(\mathcal I))\to H_*.
	\]
	Equivalently: after tree normalization, all residual transport data are
	recorded exactly by the based loop representation into the base
	stabilizer.
\end{theorem}
\begin{proof}
	We divide the proof into four steps.
	\smallskip
	\noindent
	\emph{Step 1: every gauge class admits a tree-normalized
		representative.}
	This is \cref{thm:DO-tree-normalize}. Hence every gauge-equivalence
	class contains at least one tree-normalized lift datum.
	\smallskip
	\noindent
	\emph{Step 2: every tree-normalized lift determines a homomorphism
	\(\rho:\Omega_{i_0}(\mathcal G(\mathcal I))\to H_*\).}
	Given a tree-normalized lift datum \((x,g)\), define \(\rho_g\) by
	\cref{def:DO-rho}. By \cref{lem:DO-holonomy},
	\[
		\rho_g(\ell)\in H_*
		\qquad
		\text{for all }\ell,
	\]
	and by \cref{lem:DO-rho-hom}, \(\rho_g\) is a homomorphism.
	\smallskip
	\noindent
	\emph{Step 3: every such homomorphism reconstructs a unique
		tree-normalized lift.}
	This is exactly \cref{prop:DO-reconstruct}. Therefore tree-normalized
	lifts are in bijection with homomorphisms
	\[
		\rho:\Omega_{i_0}(\mathcal G(\mathcal I))\to H_*.
	\]
	\smallskip
	\noindent
	\emph{Step 4: residual gauge acts by \(H_*\)-conjugation.}
	Let \((x,g)\) and \((x',g')\) be two tree-normalized lifts belonging to
	the same gauge-equivalence class. Since both are tree-normalized, all
	representatives are equal to \(x_*\), and all tree morphisms have
	trivial transport.
	Let \(h=(h_i)\) be a gauge modification with
	\[
		(x',g')=(x^h,g^h).
	\]
	Because
	\[
		x_i'=x_i=x_*
		\qquad
		\text{for all }i,
	\]
	one has
	\[
		h_i\cdot x_*=x_*,
	\]
	hence each \(h_i\in H_*\).
	Now let \(\tau:i\to j\) be a tree morphism. Since both lifts are
	tree-normalized,
	\[
		e=g_\tau'=h_j g_\tau h_i^{-1}=h_j e h_i^{-1},
	\]
	so
	\[
		h_j=h_i.
	\]
	As the tree is connected, it follows that all \(h_i\) are equal to a
	single element
	\[
		k\in H_*.
	\]
	By \cref{lem:DO-rho-conj},
	\[
		\rho_{g'}(\ell)=k\,\rho_g(\ell)\,k^{-1}
		\qquad
		\text{for all }\ell.
	\]
	Thus gauge-equivalent tree-normalized lifts determine \(H_*\)-conjugate
	homomorphisms.
	Conversely, let \(\rho,\rho':\Omega_{i_0}(\mathcal G(\mathcal I))\to H_*\)
	be two homomorphisms that are \(H_*\)-conjugate, say
	\[
		\rho'(\ell)=k\,\rho(\ell)\,k^{-1}
		\qquad
		(k\in H_*).
	\]
	Let \((x,g)\) and \((x',g')\) be the corresponding tree-normalized lifts
	given by \cref{prop:DO-reconstruct}. Then the constant gauge
	transformation \(h_i:=k\) transforms \((x,g)\) into \((x',g')\). Indeed,
	for any \(\alpha:i\to j\),
	\[
		g_\alpha^h
		=
		k\,g_\alpha\,k^{-1}
		=
		k\,\rho(\lambda(\alpha))\,k^{-1}
		=
		\rho'(\lambda(\alpha))
		=
		g_\alpha'.
	\]
	Since both reconstructed lifts are tree-normalized, this proves that
	they are gauge-equivalent.
	Combining the four steps yields the classification.
\end{proof}
% ============================================================
\section{Endpoint-determined transport}
% ============================================================
The obstruction representation measures exactly the failure of endpoint
determinacy. In this chapter, endpoint-determined transport means
endpoint-determined transport on the groupoid completion
\(\mathcal G(\mathcal I)\).
\begin{corollary}[Endpoint-determined transport]
	\label{cor:DO-endpoint}
	For a connected indexing category, the following are equivalent.
	\begin{enumerate}
		\item Transport on the groupoid completion \(\mathcal G(\mathcal I)\) is
		      endpoint-determined.
		\item Every loop holonomy is trivial.
		\item The obstruction representation
		      \[
			      \rho_g:\Omega_{i_0}(\mathcal G(\mathcal I))\to H_*
		      \]
		      is the trivial homomorphism.
	\end{enumerate}
\end{corollary}
\begin{proof}
	We first prove \textup{(i)}\(\Rightarrow\)\textup{(ii)}.
	Assume transport is endpoint-determined on \(\mathcal G(\mathcal I)\).
	Let
	\[
		\ell\in \Omega_{i_0}(\mathcal G(\mathcal I))
	\]
	be any loop at \(i_0\). Since \(\ell\) and \(\id_{i_0}\) have the same
	source and target, endpoint-determinacy implies
	\[
		g_\ell=g_{\id_{i_0}}=e.
	\]
	Thus every loop holonomy is trivial.
	Next, \textup{(ii)}\(\Rightarrow\)\textup{(iii)} is immediate from the
	definition of \(\rho_g\): if every loop holonomy is trivial, then
	\[
		\rho_g(\ell)=g_\ell=e
		\qquad
		\text{for every }\ell\in \Omega_{i_0}(\mathcal G(\mathcal I)),
	\]
	so \(\rho_g\) is the trivial homomorphism.
	Finally, we prove \textup{(iii)}\(\Rightarrow\)\textup{(i)}.
	Assume \(\rho_g\) is trivial. Let
	\[
		\alpha,\beta:i\to j
	\]
	be morphisms in \(\mathcal G(\mathcal I)\) with the same source and
	target. Then
	\[
		\omega:=\beta^{-1}\circ \alpha
	\]
	is a loop at \(i\). Conjugating by the tree paths gives a based loop at
	\(i_0\),
	\[
		\lambda(\omega)=\gamma_i^{-1}\circ \omega\circ \gamma_i
		\in
		\Omega_{i_0}(\mathcal G(\mathcal I)).
	\]
	Since \(\rho_g\) is trivial,
	\[
		g_{\lambda(\omega)}=e.
	\]
	By definition of \(\lambda\),
	\[
		g_{\lambda(\omega)}
		=
		g_{\gamma_i}^{-1}g_\omega g_{\gamma_i},
	\]
	hence
	\[
		g_\omega=e.
	\]
	Using \(\omega=\beta^{-1}\circ \alpha\), functoriality gives
	\[
		e=g_\omega=g_{\beta^{-1}\circ \alpha}=g_\beta^{-1}g_\alpha.
	\]
	Therefore
	\[
		g_\alpha=g_\beta.
	\]
	Thus transport on \(\mathcal G(\mathcal I)\) is endpoint-determined.
\end{proof}
\begin{corollary}[Flatness isolates object-level enrichment]
	A connected representative lift is gauge-equivalent to one with trivial
	transport if and only if its obstruction representation is trivial.
\end{corollary}
\begin{proof}
	Assume first that the obstruction representation is trivial. By
	\cref{cor:DO-endpoint}, transport is endpoint-determined on
	\(\mathcal G(\mathcal I)\). Let \((x',g')\) be a tree-normalized
	representative of the gauge class, given by
	\cref{thm:DO-tree-normalize}. Since \((x',g')\) is tree-normalized, one
	has
	\[
		g'_\tau=e
		\qquad
		\text{for every }\tau\in \Mor(\mathcal G(\mathcal T)).
	\]
	Now let \(\alpha:i\to j\) be any morphism in \(\mathcal G(\mathcal I)\).
	The tree morphism
	\[
		\tau:=\gamma_j\circ \gamma_i^{-1}:i\to j
	\]
	has the same source and target as \(\alpha\). Since transport is
	endpoint-determined,
	\[
		g'_\alpha=g'_\tau=e.
	\]
	Hence every morphism of \(\mathcal G(\mathcal I)\), and therefore every
	morphism of \(\mathcal I\), carries trivial transport in the
	tree-normalized representative. Thus the original lift is
	gauge-equivalent to one with trivial transport.
	Conversely, if the lift is gauge-equivalent to one with trivial
	transport, then in that representative
	\[
		g_\ell=e
		\qquad
		\text{for every }\ell\in \Omega_{i_0}(\mathcal G(\mathcal I)),
	\]
	so the obstruction representation is trivial. Since triviality is
	preserved under gauge conjugation, the original lift also has trivial
	obstruction representation.
\end{proof}
% ============================================================
\section{Interpretation}
% ============================================================
The classification theorem shows that the morphism locus of enrichment
has a rigid internal structure.
Tree transport carries no gauge-invariant content.
The residual content is carried entirely by the loop representation into
the stabilizer of a single representative. Thus the obstruction is not,
in general, a single group element, but a conjugacy class of
homomorphisms
\[
	\Omega_{i_0}(\mathcal G(\mathcal I))\to H_*.
\]
This is the exact descent--obstruction form of the two-locus spine:
\begin{itemize}
	\item the object locus is representative choice;
	\item the morphism locus is the obstruction representation.
\end{itemize}
No third enrichment appears, because the action groupoid contains only
objects in \(X\) and morphisms labeled by elements of \(G\).
This loop representation is the algebraic residue that persists into the
later chapters. In \cref{ch:bridge} it is identified with the finite
triangle witness on the minimal directed triangle subsystem, from there
it drives the observer-level irreversible descent of \cref{ch:flow},
continues through the stitched refinement/observer arena of
\cref{ch:stitching}, and only later does the obstruction chain pass
into graded commutator defects and curvature after smooth
realization.

\section{Conclusion}
\label{sec:descent-obstruction-conclusion}
The loop-side classification is complete: after tree normalization, all
gauge-invariant morphism enrichment is concentrated in the conjugacy class
of the obstruction representation
\(\Omega_{i_0}(\mathcal G(\mathcal I))\to H_*\).
This identifies descent failure as the exact global residue of transport
nontriviality.

Accordingly, \cref{ch:bridge} proves that this global loop
obstruction and the finite triangle witness are two presentations of
one invariant datum, and \cref{ch:flow} then derives observer-level
irreversible descent from that same static transport obstruction.

\begin{chapterclosing}
\closingitem{Established}{After tree normalization, gauge-invariant morphism enrichment is exactly the conjugacy class of the loop obstruction representation.}
\closingitem{Carried forward}{The global loop obstruction is identified next with the finite triangle witness.}
\end{chapterclosing}
